% Generated draft: requires mathematical and source review.
\documentclass[11pt]{article}
\usepackage[margin=25mm]{geometry}
\usepackage{fontspec}
\setmainfont{Latin Modern Roman}
\usepackage{amsmath,amssymb,mathtools,unicode-math}
\setmathfont{Latin Modern Math}
\usepackage{enumitem,needspace,xurl,hyperref}
\setlength{\parindent}{0pt}
\setlength{\parskip}{5pt}
\setlength{\emergencystretch}{3em}
\newcommand{\N}{\mathbb{N}}
\newcommand{\Z}{\mathbb{Z}}
\newcommand{\Q}{\mathbb{Q}}
\newcommand{\R}{\mathbb{R}}
\newcommand{\C}{\mathbb{C}}
\newcommand{\F}{\mathbb{F}}
\newcommand{\A}{\mathbb{A}}
\newcommand{\PP}{\mathbb{P}}
\newcommand{\E}{\mathbb{E}}
\newcommand{\eps}{\varepsilon}
\newcommand{\dd}{\,\mathrm{d}}
\begin{document}
\section*{UnsolvedMath ID: 10900044}
\textbf{Source problem code:} catalogue entry
\subsection{Definitions and mathematical statement}
For a Kleinian group \(\Gamma\leq\operatorname{PSL}_2(\C)\), its limit set \(\Lambda(\Gamma)\subset\partial_\infty\mathbb H^3\cong S^2\) is the set of boundary accumulation points of an orbit in \(\mathbb H^3\). Its convex core is
\[
 C_\Gamma=\operatorname{Hull}(\Lambda(\Gamma))/\Gamma,
\]
in the setting where this quotient and its boundary have the intended meaning; convex cocompactness means compactness of the convex core.

The source defines a \emph{graph Kleinian group} to be convex cocompact with the double of \(C_\Gamma\) a graph manifold. In the usual manifold interpretation the double is obtained by gluing two copies along their boundary, and a graph manifold is built from Seifert-fibered pieces along tori. The source does not supply torsion or manifold-versus-orbifold conventions for this definition.

It defines an \emph{iterated graph-Kleinian group} to be one whose Bowditch decomposition contains only hanging Fuchsian pieces and graph Kleinian pieces. The underlying Bowditch/JSJ convention, its applicability hypotheses, and the handling of elementary vertices are not specified in the excerpt; ``hanging Fuchsian'' is retained as the term for the surface-type pieces in that decomposition.

\emph{Classification problem (Walsh).} Characterize the limit sets of graph Kleinian groups and of iterated graph-Kleinian groups, with these classes kept distinct. Find intrinsic necessary and sufficient conditions, or an equivalent classification, for a set to occur as such a \(\Lambda(\Gamma)\) in the intended sense. The excerpt does not state whether limit sets are considered as abstract topological spaces, embedded subsets of \(S^2\), or up to a geometric equivalence such as quasisymmetry. Those are different classification targets and must be fixed; none is silently chosen.
\subsection{Short English statement}
Classify the limit sets of convex-cocompact Kleinian groups whose convex-core doubles are graph manifolds, and of groups assembled from hanging Fuchsian and graph Kleinian pieces in the Bowditch decomposition. The intended equivalence of limit sets and decomposition conventions are omitted.
\subsection{Sources}
\begin{itemize}
\item Local UnsolvedMath registry, permanent ID 10900044, source code catalogue entry; attribution Walsh, subsection 5.19, ``Limit sets of graph Kleinian groups.''
\item Delp, Hoffoss and Manning, \emph{Problems In Groups, Geometry, and Three-Manifolds} (2015), subsection 5.19, PDF page 5. Supplied URLs: \url{https://arxiv.org/pdf/1512.04620#page=5} and \url{https://arxiv.org/abs/1512.04620}; recorded access 2026-07-29; extraction source-tex checked against rendered PDF text. No source-side resolution is recorded.
\item The supplied summary refers to Walsh's limit-set framework, mentions ``The boundary of the convex core'' without full bibliographic details, and suggests a relationship to the Kapovich--Kleiner program and generalizations of carpets. These contextual identifications are not independently verified characterization theorems.
\item Supplied status open and OPEN-TRIAGE. The summary reports no verified complete characterization; the 2015 source presents the item as open and no complete modern primary-literature audit is recorded. Status evidence and redistribution terms for the publicly accessible author-submitted TeX and PDF remain NEEDS\_REVIEW.
\end{itemize}
\end{document}
